\section{Introduction}
\label{sec:introduction}

\subsection{The Zero-Hyperparameter Redistricting Algorithm}

Every redistricting algorithm in the B-series platform requires at
least one tunable hyperparameter.
\metis{} requires the imbalance factor \texttt{ufactor}~\citep{karypis1998}.
Simulated Annealing (B.19) requires an initial temperature factor and
step count per tract~\citep{b19paper}.
CVD (B.22) requires an iteration count~\citep{b22paper}.
\cs{} (B.16) requires a convergence threshold $T$~\citep{b16paper}.

This paper introduces \bfsg{}, a Structure-layer algorithm with no
tunable hyperparameters beyond the standard compositor inputs
(adjacency graph, population weights, balance tolerance).
The algorithm places $k$ seeds at maximally geographically spread
locations---via BFS-farthest selection from a population-weighted
starting seed---and then grows each district outward by repeatedly
assigning unassigned tracts to the adjacent district with the greatest
population deficit.
One pass through the priority queue produces a complete, contiguous
assignment.

\bfsg{} is not an optimiser.
It minimises nothing.
It does not target edge cut, does not minimise geographic distance to
centroids, and does not maximise Polsby-Popper.
It is a \emph{greedy packer}: it fills space by geographic proximity,
driven toward balance by a population-deficit priority function.
The result is the simplest geographic partition of $k$ contiguous,
population-balanced districts that can be constructed in $O(n \log n)$.

\subsection{Motivation and Use Cases}

\bfsg{} is useful in three roles.

\paragraph{Role 1: Fast initial plan (MCMC warm-start).}
MCMC redistricting chains (\recom{}, Forest ReCom) require an initial
feasible plan.
The default is a \metis{} plan.
\bfsg{} provides an alternative initial plan that is more geographically
structured---each district originates from a distinct geographic seed and
grows outward by proximity, producing districts that are roughly
geographically contiguous without the non-locality that METIS's multilevel
coarsening occasionally introduces.
Section~\ref{sec:warmstart} analyzes whether BFS-grown plans accelerate
MCMC convergence.

\paragraph{Role 2: Compactness baseline (geometric floor for EC).}
Algorithmic comparisons require a reference point: what partition quality
do you get for free, without any optimisation objective?
\bfsg{} answers this question for edge-cut quality.
Its edge-cut scores establish the \emph{EC geometric floor}---the worst-case
boundary structure produced by pure proximity-based packing.
For Polsby-Popper compactness the picture is different: \bfsg{} achieves
higher PP than \metis{} because geographic proximity packing naturally
produces rounder districts, while METIS's KL refinement optimises EC at
some PP cost.
The phrase ``geometric floor'' applies to EC, not PP.

\paragraph{Role 3: Legal argument (``simplest geographic partition'').}
In redistricting litigation, plaintiffs and defendants invoke compactness
standards.
A submitted plan can be compared not only to enacted maps and ensemble
plans but to the most primitive possible geographic partition.
If a submitted plan outperforms \bfsg{}---a single-pass, zero-optimisation
algorithm---on compactness, it demonstrates that at least some deliberate
geographic care was taken.
Conversely, if a plan \emph{underperforms} \bfsg{}, a simple greedy packer
would have done better.

\subsection{Internal Precedent: METIS Section 4}

BFS-based initial partitioning is not a new idea.
\citet{karypis1998} describe it in Section~4 of the METIS paper as the
algorithm used internally to partition the coarsened graph before
Kernighan-Lin refinement begins.
In METIS's usage, BFS partitioning is a throwaway step: its output is
immediately refined by KL.

This paper elevates METIS's internal BFS step to a standalone algorithm.
We ask: how good is BFS partitioning \emph{on its own}, without KL
refinement?
The answer---that it achieves higher EC (worse boundary quality) than
METIS but higher PP (better Polsby-Popper) than METIS, while matching
METIS in runtime---defines its operational envelope.

\subsection{Main Results}

\paragraph{Algorithm.}
We define \bfsg{} bisection (Section~\ref{sec:algorithm}) as a
standalone Structure-layer algorithm.
The algorithm has zero tunable parameters beyond the standard compositor
inputs.
A population-weighted seed selection followed by BFS-farthest spread
of remaining seeds ensures geographic diversity.
The priority function---current district population as priority (lower
is more urgent)---drives the growth toward balance.

\paragraph{Contiguity guarantee.}
We prove (Proposition~\ref{prop:contiguity}) that every district produced
by \bfsg{} is contiguous by construction.
The BFS growth frontier guarantees that each tract is added to a district
by adjacency; no disconnected assignment is possible unless the input
subgraph is itself disconnected.

\paragraph{Complexity.}
\bfsg{} runs in $O(n \log n)$ per bisection node (Proposition~\ref{prop:complexity}),
dominated by the priority queue over $n$ tracts with $O(\log n)$
per insertion.
This matches the empirical runtime of a single \metis{} call and is
substantially faster than SA (B.19) or CVD (B.22).

\paragraph{Empirical comparison.}
On NC ($k=14$), WI ($k=8$), and TX ($k=38$), \bfsg{} achieves
mean Polsby-Popper scores $5$--$12\%$ \emph{above} \metis{}
(Section~\ref{sec:comparison}), while producing $8$--$15\%$ higher edge
cut.
The geometric floor claim holds for EC: \bfsg{} produces the highest EC
of the three algorithms compared, confirming that BFS growth without
boundary refinement leaves substantial edge-cut quality on the table.
For PP-based compactness, \bfsg{} is not a lower bound; it is between
\metis{} and \cvd{}.

\paragraph{Warm-start analysis.}
Section~\ref{sec:warmstart} shows that \recom{} chains warm-started from
\bfsg{} plans have comparable or faster autocorrelation decay compared
to chains warm-started from \metis{} plans on NC and WI,
suggesting that the more geographically structured initial plan does not
impede MCMC mixing.

\subsection{Paper Structure}

Section~\ref{sec:background} reviews BFS-based partitioning in the
multilevel graph partitioning literature and the role of seed selection
in partition quality.
Section~\ref{sec:algorithm} gives the formal \bfsg{} algorithm with
pseudocode, seed derivation, and the contiguity proof.
Section~\ref{sec:comparison} presents empirical comparison against \metis{}
and \cvd{} on NC, WI, and TX.
Section~\ref{sec:warmstart} analyzes BFS Growth as an MCMC warm-start.
Section~\ref{sec:conclusion} situates \bfsg{} in the B-series algorithm
landscape.

\subsection{Relation to Prior Work}

\bfsg{} is a Structure-layer algorithm, composing cleanly with all
SeedCompositor strategies (\cs{}, \be{}) and all WeightsOverride
configurations (B.2).
It is the geometric analogue of \cvd{}~\citep{b22paper}: both use
geographic proximity as the primary signal, but \cvd{} iterates toward
centroidal Voronoi equilibria while \bfsg{} makes a single pass.
\citet{karypis1998} used BFS partitioning internally but did not
evaluate it as a standalone method.
No prior redistricting literature has examined pure BFS growth as a
standalone algorithm or established its role as the geometric
compactness floor.
